\section{Differentiating the projected Gaussian integrals}
\label{sec:gaussian-differentiation}
The Gaussian used in a local activity varies with the coarse field.
Differentiating its expectation therefore involves both the activity
and the Gaussian density. We first reduce to the coordinates read by
that activity, and then use the compensated change of variables to
put all nearby integrals on one reference probability space.

\subsection{Projection and the relative density}
Let $C$ be the positive covariance of
Section~\ref{sec:local-covariance}, in counting colour coordinates.
For a finite coordinate set $I$, let $C_I$ denote its principal
submatrix. Projection of the centered Gaussian with covariance $C$
is the centered Gaussian with covariance $C_I$. Indeed its
characteristic function at $\xi\in\mathbb R^I$ is
$\exp(-\langle\xi,C_I\xi\rangle/2)$, by extending $\xi$ by
zero. A function depending only on these coordinates can therefore
be integrated against this marginal; the other coordinates need
not be independent. The marginal density is
\begin{equation}\label{eq:gaussian-local-density}
 p_C(z)=(2\pi)^{-N/2}
   \exp\{-\tfrac12\log\det C-\tfrac12 z^TC^{-1}z\},
 \qquad N=|I|,
\end{equation}
where here and below $C$ denotes the projected covariance when
there is no ambiguity. Positivity of a principal submatrix follows
by extending a nonzero vector by zero. Formula
\eqref{eq:gaussian-local-density} follows by diagonalizing $C$ and
changing variables in the product of $N$ one-dimensional standard
Gaussian densities. It is also valid in dimension zero, with
determinant and density equal to one.

Fix a base coarse parameter $y_0$ and write $C_0=C(y_0)$. For fixed
frozen Haar variables, the compensated transformation induces a
diffeomorphism $T_y$ of the $I$ coordinates with $T_{y_0}z=z$ and
positive Jacobian $J_y(z)$. The exact expectation identity is
\begin{equation}\label{eq:gaussian-relative-identity}
 \int F_y(w)p_{C(y)}(w)\,dw
 =\int J_y(z)\frac{p_{C(y)}(T_yz)}{p_{C_0}(z)}
                  F_y(T_yz)\,d\gamma_{C_0}(z).
\end{equation}
The positive reference density cancels, leaving the ordinary
change-of-variables formula on the right. The set $I$ includes all
colours of every bond in the activity's enlarged read support.
The transformation acts bond by bond, and both it and its Jacobian
on $I$ depend only on the old coordinates in $I$. Locality therefore
identifies the transformed marginal integrand with the original
activity evaluated on the transformed full field. This explains
why the determinant in this formula is $\det C_I$, rather than the
determinant of the covariance on the full torus.

\subsection{Polynomial rank costs at the base point}
There is a useful finite-order improvement over estimating a full
determinant by its permutation expansion. If $A(y_0)=0$ is an
$N\times N$ matrix, multilinearity in the rows gives
\begin{equation}\label{eq:gaussian-small-minors}
 \det(1+A(y))=\sum_{J\subseteq\{1,\ldots,N\}}\det A_J(y).
\end{equation}
To verify the formula, in each row choose either the row of $A$
or the corresponding identity row. The identity rows force their
columns, leaving precisely the principal minor on the chosen set.
Every term of $\det A_J$ is a product of $|J|$ functions vanishing
at $y_0$. Its derivatives of order less than $|J|$ vanish, because
at least one factor receives no derivative. Thus derivatives through
order $j$ in \eqref{eq:gaussian-small-minors} use only $|J|\le j$.

Suppose every entry belongs to $\mathcal B_j(v,1;y_0)$, with
$v\ge0$. A minor of size $r\le j$ has at most $r!$ permutation
terms. The product rule bounds each term by amplitude
$(1+v)^j$ and derivative scale $j$. There are at most
$(j+1)(1+N)^j$ subsets of size at most $j$, since
$\binom Nr\le N^r$. Consequently
\[
 \det(1+A)\in
 \mathcal B_j\bigl((j+1)(1+N)^j j!(1+v)^j,j;y_0\bigr).
\]
For $j=0$ this just states the value one.

The relative density itself needs one more bookkeeping step, because
the quadratic form must be priced without an operator norm that has
not been estimated. We use the row-sum norm
\[
 \|K\|_\infty=\max_i\sum_k|K_{ik}|
\]
of an $N\times N$ matrix, which is the operator norm on
$(\mathbb R^N,\|\cdot\|_\infty)$ and is submultiplicative. It is
bounded by every exponentially weighted row price of
Section~\ref{sec:local-covariance}, since the weights are at least one.
For vectors $x,x'$ and a matrix $K$,
\begin{equation}\label{eq:gaussian-l1-linf}
 |x^TKx'|\le\sum_i|x_i|\sum_k|K_{ik}||x'_k|
 \le\|x\|_1\|K\|_\infty\|x'\|_\infty,
 \qquad |\operatorname{tr}K|\le N\|K\|_\infty .
\end{equation}
For a multilinear derivative $D^bK(y_0)$ with matrix values, we write
$\|D^bK(y_0)\|_\infty\le c$ if
$\|D^bK(y_0)[h_1,\ldots,h_b]\|_\infty\le c\|h_1\|\cdots\|h_b\|$ for all
directions $h_1,\ldots,h_b$; likewise for vector values and $\|\cdot\|_1$.

\begin{lemma}\label{lem:gaussian-base-density}
Let $N=|I|\ge1$, $j\ge0$, and let $C(y)$ be a $C^j$ family of
symmetric positive definite $N\times N$ matrices near $y_0$, with
$C_0=C(y_0)$. Let $w(y)\in\mathbb R^N$ be $C^j$ near $y_0$, with
$w(y_0)=z$. Suppose that $P,Q,v,W\ge0$ satisfy
\begin{equation}\label{eq:gaussian-base-inputs}
 \|C_0^{-1}\|_\infty\le P,\qquad
 \|D^bC(y_0)\|_\infty\le Qv^b,\qquad
 |D^rw_i(y_0)[h_1,\ldots,h_r]|\le W\|h_1\|\cdots\|h_r\|
\end{equation}
for $1\le b\le j$, $0\le r\le j$ and every $i\in I$. Put
$\kappa=\max(1,P)\max(1,Q)$. Then the relative density
\[
 R(y)=\frac{p_{C(y)}(w(y))}{p_{C_0}(z)}
\]
satisfies
\begin{equation}\label{eq:gaussian-base-density}
 R\in\mathcal B_j\bigl(1,\,K_j\,N(1+W)^2\kappa^3\max(1,v);\,y_0\bigr),
\end{equation}
where $K_j\ge1$ depends only on $j$. For $N=0$ the ratio is identically
one.
\end{lemma}

\begin{proof}
By \eqref{eq:gaussian-local-density}, $R=e^L$ with
\[
 L(y)=-\tfrac12\{\log\det C(y)-\log\det C_0\}
      -\tfrac12\{w(y)^T\Gamma(y)w(y)-z^T\Gamma(y_0)z\},
 \qquad \Gamma=C^{-1},
\]
and $L(y_0)=0$. Throughout, $c_r,c'_r,c''_r$ denote constants
depending only on $r$, and $1\le r\le j$.

\emph{Inverse.} Differentiating $\Gamma C=1$ gives
$D_h\Gamma=-\Gamma(D_hC)\Gamma$. Induction on $r$ shows that
$D^r\Gamma(y_0)[h_1,\ldots,h_r]$ is a sum of at most $c_r$ terms of the
form
\[
 \pm\Gamma_0(D^{b_1}C)\Gamma_0(D^{b_2}C)\cdots(D^{b_m}C)\Gamma_0,
 \qquad b_1+\cdots+b_m=r,\ b_i\ge1,
\]
in which the directions $h_1,\ldots,h_r$ are distributed among the
derivatives of $C$, all evaluated at $y_0$. (The induction step applies
$D_h$ either to one factor $\Gamma_0$, producing
$-\Gamma_0(D_hC)\Gamma_0$, or to one factor $D^{b}C$, raising $b$ by
one; this bounds $c_{r+1}$ by $(2r+1)c_r$.) By submultiplicativity and
\eqref{eq:gaussian-base-inputs}, each term has row-sum norm at most
$P^{m+1}Q^mv^r\le\kappa^{r+1}v^r$, since $m\le r$. Hence
$\|D^r\Gamma(y_0)\|_\infty\le c_r\kappa^{r+1}v^r$, and
$\|\Gamma(y_0)\|_\infty\le\kappa$.

\emph{Determinant.} Near $y_0$,
$D_h\log\det C=\operatorname{tr}(\Gamma\,D_hC)$. Differentiating this
$r-1$ more times by Leibniz's rule gives at most $2^{r-1}$ terms
$\operatorname{tr}\bigl(D^a\Gamma\cdot D^{b+1}C\bigr)$ with $a+b=r-1$.
By \eqref{eq:gaussian-l1-linf} and the inverse bound, each is at most
$N\cdot c_a\kappa^{a+1}v^a\cdot Qv^{b+1}\le c_aN\kappa^{r+1}v^r$.
Thus
\[
 |D^r\log\det C(y_0)|\le c'_rN\kappa^{r+1}v^r .
\]

\emph{Quadratic form.} Leibniz's rule for the trilinear expression
$w^T\Gamma w$ writes $D^r(w^T\Gamma w)(y_0)$ as a sum of at most $3^r$
terms $(D^aw)^T(D^b\Gamma)(D^cw)$ with $a+b+c=r$, all at $y_0$.
By \eqref{eq:gaussian-base-inputs}, $\|D^aw(y_0)\|_\infty\le W$ and
$\|D^aw(y_0)\|_1\le NW$, including $a=0$, where $D^0w(y_0)=z$. Now
\eqref{eq:gaussian-l1-linf} and the inverse bound give, for each term,
the bound $NW\cdot c_b\kappa^{b+1}\max(1,v)^b\cdot W$. Hence
\[
 |D^r(w^T\Gamma w)(y_0)|\le c''_rN W^2\kappa^{r+1}\max(1,v)^r .
\]
The Euclidean norm of $z$ and the operator norm of $\Gamma$ on
$\ell^2$ never enter.

\emph{The exponent.} Combining the two bounds, with $Q\le\kappa$ and
$W^2\le(1+W)^2$,
\[
 |D^rL(y_0)|\le K'_rN(1+W)^2\kappa^{r+1}\max(1,v)^r,\qquad
 K'_r=\tfrac12(c'_r+c''_r).
\]
Let $K=\max_{1\le r\le j}\max(1,K'_r)$ and
$S=KN(1+W)^2\kappa^3\max(1,v)$. Every factor in $S$ is at least one,
because $N\ge1$. Therefore, for $1\le r\le j$, using $\kappa^{r+1}\le
\kappa^{3r}$ and $N(1+W)^2\le\{N(1+W)^2\}^r$,
\[
 |D^rL(y_0)|\le S^r .
\]
This absorbs the constants $K'_r$ into the scale.

\emph{Exponentiation.} By the exponential partition formula used in
Section~\ref{sec:packet-estimates}, $D^re^L(y_0)$ is the sum over the
set partitions of the $r$ directions of $e^{L(y_0)}$ times the product,
over the blocks, of the derivatives of $L$ of the block sizes. Since
$L(y_0)=0$, each such product is at most $S^r$, and there are at most
$r!$ partitions (the Bell number is at most $r!$). Hence
$|D^rR(y_0)|\le r!S^r\le(j!S)^r$ for $1\le r\le j$, while $R(y_0)=1$.
This is \eqref{eq:gaussian-base-density} with $K_j=\max(1,j!)K$.
\end{proof}

In the application $w=T_yz$, the coordinate bound $W$ is that of
Section~\ref{sec:gaussian-moment-prices}, of the form
$W\le\sqrt\beta\{u(x)+2m(x)\}+F_I(z)$, and the quantities $P,Q,v$ are
the weighted row prices of the principal covariance, its inverse and
their derivatives from Sections~\ref{sec:local-covariance}
and~\ref{sec:normalization-jets}. Since $1\le F_I(z)$ and $N\le F_I(z)$,
and $1+a+F\le2(1+a)F$ for $a\ge0$, $F\ge1$, the scale in
\eqref{eq:gaussian-base-density} is at most $D'(x)F_I(z)^3$ with
$D'=4K_j\{1+\sqrt\beta(u+2m)\}^2\kappa^3\max(1,v)$, which has
upper exponent zero when $\kappa$ and $v$ do.
This is the form \eqref{eq:moment-ratio-jets} used there, before the
Jacobian factor is multiplied in. This is a finite-rank derivative
estimate at the base point. It does not yet justify taking derivatives
under an integral over unbounded $z$; that is the purpose of the next
subsection.

\subsection{An integrable bound on one open neighbourhood}
We prove the additional assertion needed for that interchange.
Let $y$ range over a finite-dimensional real parameter space and
let $C(y)$ be a $C^j$ family of symmetric matrices near $y_0$,
with $C_0$ positive definite. Suppose
\[
 C_0^{-1}\ge\lambda 1,\qquad\lambda>0.
\]
Let $z$ be an arbitrary sample and let $w_z(y)$ be a family of
$C^j$ maps such that, on one neighbourhood of $y_0$, uniformly in
$z$ and in any auxiliary frozen variables,
\begin{equation}\label{eq:gaussian-bounded-displacement}
 w_z(y)-z\in\mathcal B_j(B,1;y),\qquad B\ge0.
\end{equation}
Then there is a nonnegative function $M$, integrable for
$\gamma_{C_0}$, which bounds every derivative through order $j$ of
$p_{C(y)}(w_z(y))/p_{C_0}(z)$ on a possibly smaller common
neighbourhood.

By continuity and finite dimension, shrink the neighbourhood so
that $C(y)$ is positive definite and
$C(y)^{-1}\ge(\lambda/2)1$. For an explicit verification of the
last inequality, make every entry of $C(y)^{-1}-C_0^{-1}$ at most
$\lambda/(2(1+N))$ in absolute value. Its row and column sums are
at most $\lambda/2$, so its quadratic form is bounded in absolute
value by $(\lambda/2)\|z\|^2$. On a smaller closed neighbourhood,
all required derivatives of $C^{-1}$ and $\log\det C$ have finite
bounds. These bounds may depend on $N$ and the fixed period.

The derivatives of $w_z$ have amplitude at most $B+\|z\|$.
Thus the logarithm of the numerator has derivatives bounded by
$K(1+\|z\|)^2$ for a fixed $K$. Moreover
\[
 \|z\|^2\le2\|w_z(y)\|^2+2B^2,
\]
and hence
\[
 -\tfrac12\log\det C(y)
       -\tfrac12 w_z(y)^TC(y)^{-1}w_z(y)
 \le K_0+\lambda B^2/4-\lambda\|z\|^2/8.
\]
The repeated exponential chain rule now bounds every derivative
through order $j$ of the relative density by
\begin{equation}\label{eq:gaussian-open-envelope}
 M(z)=K_1\frac{e^{-\lambda\|z\|^2/8}
                         (1+\|z\|)^{2j}}{p_{C_0}(z)}.
\end{equation}
This is integrable against $\gamma_{C_0}$: multiplication by its
density cancels the denominator, and a polynomial times a strictly
decaying Gaussian is integrable in finite dimension. Explicitly,
$(1+t)^{2j}\le e^{2jt}$ for $t\ge0$, and completing the square
absorbs $2jt$ into half the negative quadratic exponent. The
remaining integral factors into $N$ finite one-dimensional Gaussian
integrals. The same proof works when $z$ is a measurable function
of a larger probability space with law $\gamma_{C_0}$.

\subsection{Why the activity satisfies the neighbourhood hypotheses}
For packet factors we use the locally smooth stock extension of
Section~\ref{sec:packet-estimates}, with its fixed seed-eligibility
indicator. Under the stated healthy-footprint and sector conditions,
its value at the base field equals the physical packet. We do not
differentiate the physical packet's varying flag test.

For the compensated map of
Section~\ref{sec:compensated-variation}, write
$u_0=\nu_0+\beta^{-1/2}z$ and
$F_M(u)=u+\psi(M,u)$. Its displacement in Gaussian coordinates is
\begin{equation}\label{eq:gaussian-compensated-displacement}
 T_yz-z=\sqrt\beta\,
   \{\psi(M(y),u_0)+\nu_0-\nu(y)\}.
\end{equation}
Here the fixed coupling satisfies $\beta>0$. The perturbation
$\psi$ and all its parameter derivatives are uniformly bounded
for all $u_0$: outside the fixed radius-$2d_*$ ball of
Section~\ref{sec:compensated-variation} it is zero, and
inside that ball compactness applies. The unmasked mean and its
derivatives are bounded on a common neighbourhood uniformly in
the frozen Haar variables. Indeed the frozen cut logarithms are
bounded smooth functions of the coarse variation; the finitely
many source sums, row matrices and principal inverses have uniform
derivative bounds there. All principal determinants stay nonzero
after one common shrink. Formula
\eqref{eq:gaussian-compensated-displacement} proves
\eqref{eq:gaussian-bounded-displacement}. The same compact argument
and the finite product rule bound the Jacobian derivatives uniformly.

The remaining activity integrand also has bounded derivatives on
one neighbourhood, uniformly in the Gaussian sample. The required
localization is worth spelling out. Compare two fluctuation arrays
$u,\widetilde u$ such that at each bond either they agree or both
have norm greater than $4a\tau$. Every cutoff of width at most
$2a\tau$ is then equal on the two arrays, and if that cutoff is
nonzero the arrays agree at that bond. A masked chart logarithm
therefore agrees as well. A packet's nonzero protected product
forces agreement throughout its $304B$ region, which contains the
read stencils of all its $200B$ action terms. Thus its logarithmic
exponent and its complete value agree.

For a linear Gaussian ratio term, the score vanishes on inactive
bonds and its hot endpoint carries the doubled cutoff. For a
quadratic ratio term, every nonzero matrix itinerary has two hot
endpoints. To check this, $AC-1$ is zero if either coordinate is
inactive, since both $A$ and $C$ equal the identity on that
coordinate. In a product of copies of $-(AC-1)$ followed by $A$,
the first nonzero factor forces the initial endpoint hot, and the
last factor cannot join a hot coordinate to an inactive one.
Both endpoint cutoffs therefore force equality of the coordinates
actually used by the quadratic term. This argument holds nearby
as long as the principal inverses exist; it does not require the
nearby point to satisfy the original sector inequalities.

It follows that every fixed activity factor is unchanged when
large initial fluctuation coordinates are replaced by bounded
ones, provided this comparison is made through the compensated
map and the fine field is held fixed. Apply radial truncation at
radius $3d_*$ to the initial fluctuation coordinates
$u_0=\nu_0+\beta^{-1/2}z$ of each bond. These do not depend on the
parameter $y$, so the lack of smoothness of radial truncation in its
argument plays no role. The truncation leaves the closed
radius-$3d_*$ ball unchanged, and in particular the radius-$2d_*$
ball, which is all that is used. Outside that ball $F_M$ is the
identity; both the old and truncated coordinates then remain beyond
the cutoff threshold, since $8a\tau\le d_*/2$. Inside it, the
coordinates agree. The resulting
bounded input family has uniform parameter derivatives by the
compact chart construction. The frozen fine factors remain merely
bounded measurable functions: their regular activities, forced
activities, likelihoods and contractions use the supplied
finite-volume bounds for the chosen labels. No differentiation
of these frozen factors is required.

For each fixed finite product of activity factors, finite sums,
products and smooth compositions on bounded sets now give a
common bound on all derivatives through order $j$. The local
action, normalization and chart functions used here are the
unmasked smooth extensions already constructed. Hard selection
indicators are fixed. Together with
\eqref{eq:gaussian-open-envelope} and the Jacobian bound, this
gives an integrable envelope for the entire right-hand side of
\eqref{eq:gaussian-relative-identity}, also after integrating the
frozen variables against their Haar probability law.

\subsection{Interchange and quantitative use}
Here is the precise final differentiation argument. Suppose
$G(y,\omega)$ is measurable in $\omega$, $C^j$ in $y$ on an open
neighbourhood, and for each $r\le j$ its $r$th derivative is
bounded there by an integrable function $M_r(\omega)$.
Then
\begin{equation}\label{eq:gaussian-differentiate-integral}
 D^r\int G(y,\omega)\,d\mu(\omega)
       =\int D_y^rG(y,\omega)\,d\mu(\omega),\qquad r\le j.
\end{equation}
Derivative measurability follows in a fixed finite basis by taking
limits of difference quotients of measurable functions. Dominated
convergence first gives continuity of each derivative integral.
For the induction step, the fundamental theorem of calculus on a
short parameter segment writes its difference quotient as the
integral of the next derivative along that segment. The bound
$M_{r+1}$ permits dominated convergence as the segment shrinks.
Starting at order zero proves the derivative identity and its
continuity successively through order $j$.

Apply this to the common reference expression above. It proves
smoothness of each projected activity integral. If a sharper
estimate at the base point gives
$G(\cdot,\omega)\in\mathcal B_j(q(\omega),D;y_0)$ and
$\int q\,d\mu\le Q$, then
\eqref{eq:gaussian-differentiate-integral} and the triangle
inequality give
\[
 \int G(\cdot,\omega)\,d\mu(\omega)
            \in\mathcal B_j(Q,D;y_0).
\]
The neighbourhood envelope establishes the validity of
differentiation; the sharper moment estimate supplies its
period-independent price. Those moment estimates must be paid
separately. In particular the period-dependent constants used to
justify \eqref{eq:gaussian-differentiate-integral} are never used
as uniform renormalization bounds.
